\begin{lexlean}{Guide}
\useglossary{lexlean.std.bool@1.2.0}
\useglossary{lexlean.std.nat@1.2.0}
\useglossary{prism.arch@1.0.0}
\importmodule{Foundation.View.Text.V3.Model}
\title{Prism-code workbench view}
\begin{semanticmodule}
definition{
  name: purpose,
  parameters: [],
  result: String,
  body: "Checks a Prism dataflow project script against its declared memory and PII-encryption policies and prints the canonical project."
}

definition{name: checkingText, parameters: [], result: String, body: "Checking the project script."}

definition{
  name: acceptedText,
  parameters: [],
  result: String,
  body: "Accepted: the project meets its policies. The canonical project follows."
}

definition{
  name: rejectedText,
  parameters: [],
  result: String,
  body: "Rejected: the project breaks a policy. Each finding names the node and the policy it breaks; the canonical project follows."
}

definition{
  name: errorText,
  parameters: [],
  result: String,
  body: "Error: the script could not be checked. The report names the first failing line and its code; the fix for that code is listed under error and finding codes."
}

definition{
  name: overflowText,
  parameters: [],
  result: String,
  body: "The report is too large for the page. Check this script from the command line."
}

definition{
  name: sections,
  parameters: [],
  result: (List Foundation.View.Text.V3.Model::Section),
  body: (cons (record Foundation.View.Text.V3.Model::Section {id := "node", title := "node NAME KIND MEMORY", paragraphs := (cons "Declares a node; KIND is input, operator, encrypt or output and MEMORY is the bytes it requires." (nil String))}) (cons (record Foundation.View.Text.V3.Model::Section {id := "edge", title := "edge SOURCE TARGET", paragraphs := (cons "Adds a dataflow edge between two nodes declared on earlier lines." (nil String))}) (cons (record Foundation.View.Text.V3.Model::Section {id := "pii", title := "pii NAME", paragraphs := (cons "Marks an input node as carrying personal data." (nil String))}) (cons (record Foundation.View.Text.V3.Model::Section {id := "memory", title := "policy memory LIMIT", paragraphs := (cons "No node may require more than LIMIT bytes." (nil String))}) (cons (record Foundation.View.Text.V3.Model::Section {id := "encrypt", title := "policy encrypt-pii", paragraphs := (cons "Personal data must pass an encrypt node before it reaches any output." (nil String))}) (nil Foundation.View.Text.V3.Model::Section))))))
}

definition{
  name: errors,
  parameters: [],
  result: (List Foundation.View.Text.V3.Model::ErrorHint),
  body: (cons (record Foundation.View.Text.V3.Model::ErrorHint {code := "unknown-command", title := "unknown-command", explanation := "The line starts with a word that is not a statement, or policy names an unknown policy.", recovery := "Start the line with node, edge, pii or policy memory/encrypt-pii, or delete the line.", fields := (nil String)}) (cons (record Foundation.View.Text.V3.Model::ErrorHint {code := "wrong-arity", title := "wrong-arity", explanation := "The statement has too few or too many fields.", recovery := "Give the statement exactly the fields shown for it in the statement reference.", fields := (nil String)}) (cons (record Foundation.View.Text.V3.Model::ErrorHint {code := "empty-field", title := "empty-field", explanation := "The line has two spaces in a row, which leaves a field empty.", recovery := "Separate fields with exactly one space and remove leading or trailing spaces.", fields := (nil String)}) (cons (record Foundation.View.Text.V3.Model::ErrorHint {code := "unknown-kind", title := "unknown-kind", explanation := "The node kind is not one of the four kinds.", recovery := "Use input, operator, encrypt or output as the node kind.", fields := (nil String)}) (cons (record Foundation.View.Text.V3.Model::ErrorHint {code := "invalid-number", title := "invalid-number", explanation := "A memory value is not a whole number of bytes.", recovery := "Write the memory or limit as digits only, with no sign, unit or separator.", fields := (nil String)}) (cons (record Foundation.View.Text.V3.Model::ErrorHint {code := "duplicate-node", title := "duplicate-node", explanation := "A node with this name is already declared.", recovery := "Give each node a unique name, or remove the repeated node line.", fields := (nil String)}) (cons (record Foundation.View.Text.V3.Model::ErrorHint {code := "unknown-node", title := "unknown-node", explanation := "An edge or pii line names a node that has not been declared on an earlier line.", recovery := "Declare the node with a node line before the line that uses it, or correct the name.", fields := (nil String)}) (cons (record Foundation.View.Text.V3.Model::ErrorHint {code := "self-edge", title := "self-edge", explanation := "The edge connects a node to itself.", recovery := "Connect two different nodes, or remove the edge.", fields := (nil String)}) (cons (record Foundation.View.Text.V3.Model::ErrorHint {code := "duplicate-edge", title := "duplicate-edge", explanation := "This edge is already declared.", recovery := "Remove the repeated edge line.", fields := (nil String)}) (cons (record Foundation.View.Text.V3.Model::ErrorHint {code := "output-source", title := "output-source", explanation := "The edge starts at an output node, and data cannot flow out of an output.", recovery := "Start the edge at an input, operator or encrypt node.", fields := (nil String)}) (cons (record Foundation.View.Text.V3.Model::ErrorHint {code := "input-target", title := "input-target", explanation := "The edge ends at an input node, and data cannot flow into an input.", recovery := "End the edge at an operator, encrypt or output node.", fields := (nil String)}) (cons (record Foundation.View.Text.V3.Model::ErrorHint {code := "not-an-input", title := "not-an-input", explanation := "Only an input node can be marked as carrying personal data.", recovery := "Mark the input node where the personal data enters, not a later node.", fields := (nil String)}) (cons (record Foundation.View.Text.V3.Model::ErrorHint {code := "duplicate-pii", title := "duplicate-pii", explanation := "This node is already marked as carrying personal data.", recovery := "Remove the repeated pii line.", fields := (nil String)}) (cons (record Foundation.View.Text.V3.Model::ErrorHint {code := "duplicate-policy", title := "duplicate-policy", explanation := "This policy is already stated; each policy may be stated once.", recovery := "Remove the repeated policy line, or keep the one you want.", fields := (nil String)}) (cons (record Foundation.View.Text.V3.Model::ErrorHint {code := "line-overflow", title := "line-overflow", explanation := "The script has more lines than the line counter can hold.", recovery := "Split the script into smaller projects.", fields := (nil String)}) (cons (record Foundation.View.Text.V3.Model::ErrorHint {code := "empty", title := "empty", explanation := "The request is empty.", recovery := "Enter the header dataflow-policy/1 on the first line, then the project statements.", fields := (nil String)}) (cons (record Foundation.View.Text.V3.Model::ErrorHint {code := "too-large", title := "too-large", explanation := "The request is longer than 4096 bytes.", recovery := "Shorten the script to at most 4096 bytes.", fields := (nil String)}) (cons (record Foundation.View.Text.V3.Model::ErrorHint {code := "invalid-utf8", title := "invalid-utf8", explanation := "The request is not valid UTF-8 text.", recovery := "Paste the script as plain UTF-8 text.", fields := (nil String)}) (cons (record Foundation.View.Text.V3.Model::ErrorHint {code := "missing-header", title := "missing-header", explanation := "The first line is not the header dataflow-policy/1.", recovery := "Start the script with the line dataflow-policy/1.", fields := (nil String)}) (cons (record Foundation.View.Text.V3.Model::ErrorHint {code := "memory", title := "memory", explanation := "A node requires more memory than the policy memory limit allows.", recovery := "Raise the policy memory limit, or lower the memory declared for the node the finding names.", fields := (nil String)}) (cons (record Foundation.View.Text.V3.Model::ErrorHint {code := "plaintext-pii", title := "plaintext-pii", explanation := "Personal data reaches an output without passing an encrypt node.", recovery := "Route the data through an encrypt node before the output the finding names, and remove any edge that bypasses it.", fields := (nil String)}) (nil Foundation.View.Text.V3.Model::ErrorHint))))))))))))))))))))))
}

definition{
  name: examples,
  parameters: [],
  result: (List Foundation.View.Text.V3.Model::Example),
  body: (cons (record Foundation.View.Text.V3.Model::Example {id := "accepted", label := "Compliant pipeline", description := "A camera's personal data passes an encrypt node before it reaches the screen.", values := (cons "dataflow-policy/1\npolicy memory 4000000\npolicy encrypt-pii\nnode camera input 0\nnode detector operator 2000000\nnode vault encrypt 1000\nnode screen output 0\npii camera\nedge camera detector\nedge detector vault\nedge vault screen" (nil String)), expectedResponse := "{\"version\":1,\"status\":\"ok\",\"operation\":\"check\",\"message\":\"Accepted: the project meets its policies. The canonical project follows.\",\"rows\":11,\"columns\":1,\"cells\":[\"dataflow-policy/1\",\"policy memory 4000000\",\"policy encrypt-pii\",\"node camera input 0\",\"node detector operator 2000000\",\"node screen output 0\",\"node vault encrypt 1000\",\"pii camera\",\"edge camera detector\",\"edge detector vault\",\"edge vault screen\"],\"details\":[],\"request\":\"dataflow-policy/1\\npolicy memory 4000000\\npolicy encrypt-pii\\nnode camera input 0\\nnode detector operator 2000000\\nnode vault encrypt 1000\\nnode screen output 0\\npii camera\\nedge camera detector\\nedge detector vault\\nedge vault screen\",\"response\":\"status accepted\\ndataflow-policy/1\\npolicy memory 4000000\\npolicy encrypt-pii\\nnode camera input 0\\nnode detector operator 2000000\\nnode screen output 0\\nnode vault encrypt 1000\\npii camera\\nedge camera detector\\nedge detector vault\\nedge vault screen\",\"error\":\"\"}", explanation := "The report starts with status accepted, then the canonical project."}) (cons (record Foundation.View.Text.V3.Model::Example {id := "rejected", label := "Bypassed encryption", description := "An extra edge sends detector output to the screen unencrypted, and the detector needs more memory than the limit.", values := (cons "dataflow-policy/1\npolicy memory 1000000\npolicy encrypt-pii\nnode camera input 0\nnode detector operator 2000000\nnode vault encrypt 1000\nnode screen output 0\npii camera\nedge camera detector\nedge detector vault\nedge vault screen\nedge detector screen" (nil String)), expectedResponse := "{\"version\":1,\"status\":\"ok\",\"operation\":\"check\",\"message\":\"Rejected: the project breaks a policy. Each finding names the node and the policy it breaks; the canonical project follows.\",\"rows\":14,\"columns\":1,\"cells\":[\"finding memory detector 2000000 exceeds 1000000\",\"finding plaintext-pii screen\",\"dataflow-policy/1\",\"policy memory 1000000\",\"policy encrypt-pii\",\"node camera input 0\",\"node detector operator 2000000\",\"node screen output 0\",\"node vault encrypt 1000\",\"pii camera\",\"edge camera detector\",\"edge detector screen\",\"edge detector vault\",\"edge vault screen\"],\"details\":[],\"request\":\"dataflow-policy/1\\npolicy memory 1000000\\npolicy encrypt-pii\\nnode camera input 0\\nnode detector operator 2000000\\nnode vault encrypt 1000\\nnode screen output 0\\npii camera\\nedge camera detector\\nedge detector vault\\nedge vault screen\\nedge detector screen\",\"response\":\"status rejected\\nfinding memory detector 2000000 exceeds 1000000\\nfinding plaintext-pii screen\\ndataflow-policy/1\\npolicy memory 1000000\\npolicy encrypt-pii\\nnode camera input 0\\nnode detector operator 2000000\\nnode screen output 0\\nnode vault encrypt 1000\\npii camera\\nedge camera detector\\nedge detector screen\\nedge detector vault\\nedge vault screen\",\"error\":\"\"}", explanation := "The report starts with status rejected, lists both findings, then the canonical project."}) (cons (record Foundation.View.Text.V3.Model::Example {id := "error", label := "Unknown node", description := "An edge names a node that was never declared.", values := (cons "dataflow-policy/1\nnode camera input 0\nedge camera detector" (nil String)), expectedResponse := "{\"version\":1,\"status\":\"error\",\"operation\":\"check\",\"message\":\"Error: the script could not be checked. The report names the first failing line and its code; the fix for that code is listed under error and finding codes.\",\"rows\":0,\"columns\":0,\"cells\":[],\"details\":[\"line 3 unknown-node\"],\"request\":\"dataflow-policy/1\\nnode camera input 0\\nedge camera detector\",\"response\":\"status error\\nline 3 unknown-node\",\"error\":\"unknown-node\"}", explanation := "The report starts with status error and names line 3 and the code unknown-node."}) (nil Foundation.View.Text.V3.Model::Example))))
}

definition{
  name: view,
  parameters: [],
  result: Foundation.View.Text.V3.Model::WorkbenchView,
  body: (record Foundation.View.Text.V3.Model::WorkbenchView {title := "Dataflow Policy", heading := "Dataflow Policy", purpose := (purpose), audience := "People who write Prism dataflow project scripts.", guidance := (cons "Start the script with the header dataflow-policy/1." (cons "Write one statement per line; the statement reference lists all five forms." (cons "Check the project and read the status line first." (cons "Every error and finding code below has an explanation and a fix." (nil String))))), initialOperation := "check", initialExample := "accepted", operations := (cons (record Foundation.View.Text.V3.Model::Operation {id := "check", label := "Check project", summary := "Check a project script against its memory and PII-encryption policies.", help := "Each finding names the node and the policy it breaks.", requestPrefix := "ui-check", fields := (cons (record Foundation.View.Text.V3.Model::Field {id := "script", label := "Project script", kind := (ctor Foundation.View.Text.V3.Model::FieldKind.Multiline), help := "The header dataflow-policy/1, then one statement per line; at most 4087 bytes on this page.", placeholder := "dataflow-policy/1", initialValue := "dataflow-policy/1", required := true, choices := (nil Foundation.View.Text.V3.Model::Choice)}) (nil Foundation.View.Text.V3.Model::Field)), examples := (examples)}) (nil Foundation.View.Text.V3.Model::Operation)), sections := (sections), links := (cons (record Foundation.View.Text.V3.Model::Link {label := "Source", httpsUrl := "https://github.com/UOR-Foundation/dataflow-policy"}) (nil Foundation.View.Text.V3.Model::Link)), errors := (errors), statuses := (record Foundation.View.Text.V3.Model::StatusCopy {scriptRequired := "Enable JavaScript to check a project.", initializing := "Loading the checker.", ready := "Ready.", running := (checkingText), cancelling := "Cancelling the check.", cancelled := "The check was cancelled.", timeout := "The check took too long.", loadError := "The checker could not be loaded.", computeError := "The check failed.", inputError := "Enter UTF-8 text of at most 4087 bytes.", responseError := "The report could not be displayed."}), design := (record Foundation.View.Text.V3.Model::Design {palette := (ctor Foundation.View.Text.V3.Model::Palette.Indigo), density := (ctor Foundation.View.Text.V3.Model::Density.Comfortable), width := (ctor Foundation.View.Text.V3.Model::Width.Balanced), colorScheme := (ctor Foundation.View.Text.V3.Model::ColorScheme.System)}), raw := (record Foundation.View.Text.V3.Model::RawEditor {requestPrefix := "ui-raw", label := "Raw request", help := "Send an exact dataflow-policy/1 request, header included.", placeholder := "dataflow-policy/1", initialValue := "dataflow-policy/1", submitLabel := "Check request", responseLabel := "Report", expanded := false}), response := (record Foundation.View.Text.V3.Model::DisplaySpec {title := "Report", emptyText := "No report yet.", rawLabel := "Exact report", maximumCells := (4096 : UInt32), maximumDetails := (128 : UInt32)}), execution := (record Foundation.View.Text.V3.Model::Execution {timeoutMilliseconds := (30000 : UInt32), cancelLabel := "Cancel", retryLabel := "Try again", wasmCachePolicy := (ctor Foundation.View.Text.V3.Model::WasmCachePolicy.NoStore)}), labels := (record Foundation.View.Text.V3.Model::WorkbenchLabels {navigation := "Sections", operation := "Operation", examples := "Examples", loadExample := "Load example", compute := "Check project", guide := "Statement reference", links := "Links", fields := "Fields", expectedResponse := "Expected report", details := "Details", resultTable := "Report lines", row := "Line", column := "Column", request := "Request", reset := "Reset", skipToWorkbench := "Skip to the checker"})})
}
\end{semanticmodule}
\end{lexlean}
